% !PRINTED 290
\weberpage{290}
%<<CONT
\begin{center}
\S.\,95.

\gesperrt{Relation of the characteristic to the
points of intersection.}
\end{center}

One must now observe that the designation of the points of intersection
of two curves as points of entry or of exit in \S.\,93
is essentially different from that in \S.\,94. There only
the relation of the points to the two intersecting
curves came into consideration, while in \S.\,94 the relation to
a third curve was also in question; this is completely expressed in the chosen
designations $E(\varphi, \psi)$, $A(\varphi, \psi)$ and

\weberfigureimage{weber-b1-fig10.png}{0.48}

$E(f; \varphi, \psi)$, $A(f; \varphi, \psi)$. But now we must
investigate more precisely how the two modes of designation
are related to one another.

For the sake of clarity we choose a somewhat simpler
figure than above, which at the same time illustrates all possible
circumstances (Fig.~10).

At a point $E(\varphi, \psi)$ the $\varphi$-curve enters the region
of the curve $\psi$, i.\,e.\ $\psi$ passes from positive to negative
values. If at the same time $f$ is positive, then this point is an
$E(\varphi; \psi, f)$, like the point 1 in our figure; but if $f$ is negative,
then it is an $A(\varphi; \psi, f)$, like the point 3. Thus we have:

\[
\begin{array}{ll}
\text{A point } E(\varphi, \psi) \text{ is an} & \\
\quad E(\varphi; \psi, f), \text{ if } f > 0, & \text{(point 1 of Fig. 10)} \\
\quad A(\varphi; \psi, f), \text{ if } f < 0, & \text{(point 3 of Fig. 10).} \\
\text{A point } A(\varphi, \psi) \text{ is an} & \\
\quad E(\varphi; \psi, f), \text{ if } f < 0, & \text{(point 2 of Fig. 10)} \\
\quad A(\varphi; \psi, f), \text{ if } f > 0, & \text{(point 4 of Fig. 10).}
\end{array}
\]

Now at the points $E(\varphi, \psi)$ the functional determinant
$[\varphi, \psi]$ is negative, and at $A(\varphi, \psi)$ positive (\S.\,93), and we can
therefore summarise what has been proved as follows:

$E(\varphi; \psi, f)$ is the number of points of intersection of $\varphi$, $\psi$ at
which $[\varphi, \psi]\, f < 0$,

$A(\varphi; \psi, f)$ the number of points of intersection of $\varphi$, $\psi$ at
which $[\varphi, \psi]\, f > 0$;
%CONT>>
